\documentclass[10pt,a4paper]{amsart}
\usepackage{amsmath,amssymb,amsthm,graphicx}
% For dealing with references we use the comment environment
\usepackage{verbatim}
\newenvironment{reference}{\comment}{\endcomment}
%\newenvironment{reference}{}{}
\newenvironment{slogan}{\comment}{\endcomment}
\newenvironment{history}{\comment}{\endcomment}

% For commutative diagrams we use Xy-pic
\usepackage[all]{xy}
% Theorem environments.
%
\theoremstyle{plain}
\newtheorem{theorem}[subsection]{Theorem}
\newtheorem{proposition}[subsection]{Proposition}
\newtheorem{lemma}[subsection]{Lemma}

\theoremstyle{definition}
\newtheorem{definition}[subsection]{Definition}
\newtheorem{example}[subsection]{Example}
\newtheorem{exercise}[subsection]{Exercise}
\newtheorem{situation}[subsection]{Situation}

\theoremstyle{remark}
\newtheorem{remark}[subsection]{Remark}
\newtheorem{remarks}[subsection]{Remarks}

\numberwithin{equation}{subsection}

% Macros
%
\def\lim{\mathop{\mathrm{lim}}\nolimits}
\def\colim{\mathop{\mathrm{colim}}\nolimits}
\def\Spec{\mathop{\mathrm{Spec}}}
\def\Hom{\mathop{\mathrm{Hom}}\nolimits}
\def\Ext{\mathop{\mathrm{Ext}}\nolimits}
\def\SheafHom{\mathop{\mathcal{H}\!\mathit{om}}\nolimits}
\def\SheafExt{\mathop{\mathcal{E}\!\mathit{xt}}\nolimits}
\def\Sch{\mathit{Sch}}
\def\Mor{\mathop{\mathrm{Mor}}\nolimits}
\def\Ob{\mathop{\mathrm{Ob}}\nolimits}
\def\Sh{\mathop{\mathit{Sh}}\nolimits}
\def\NL{\mathop{N\!L}\nolimits}
\def\CH{\mathop{\mathrm{CH}}\nolimits}
\def\proetale{{pro\text{-}\acute{e}tale}}
\def\etale{{\acute{e}tale}}
\def\QCoh{\mathit{QCoh}}
\def\Ker{\mathop{\mathrm{Ker}}}
\def\Im{\mathop{\mathrm{Im}}}
\def\Coker{\mathop{\mathrm{Coker}}}
\def\Coim{\mathop{\mathrm{Coim}}}

% Boxtimes
%
\DeclareMathSymbol{\boxtimes}{\mathbin}{AMSa}{"02}

%
% Macros for moduli stacks/spaces
%
\def\QCohstack{\mathcal{QC}\!\mathit{oh}}
\def\Cohstack{\mathcal{C}\!\mathit{oh}}
\usepackage{cleveref}
\crefname{originallemma}{original lemma}{original lemmas}
\Crefname{originallemma}{Original lemma}{Original lemmas}
\makeatletter
\newlabel{categories-lemma-what-gets-inverted}{{05Q9}{source}{}{sourcepointer-1}{}}
\newlabel{categories-lemma-what-gets-inverted@cref}{{[originallemma][0][]05Q9}{[1][0][]source}{}{}{}}
\makeatother
\crefname{originaldefinition}{original definition}{original definitions}
\Crefname{originaldefinition}{Original definition}{Original definitions}
\makeatletter
\newlabel{categories-definition-directed}{{002V}{source}{}{sourcepointer-2}{}}
\newlabel{categories-definition-directed@cref}{{[originaldefinition][0][]002V}{[1][0][]source}{}{}{}}
\makeatother
\crefname{originaldefinition}{original definition}{original definitions}
\Crefname{originaldefinition}{Original definition}{Original definitions}
\makeatletter
\newlabel{categories-definition-cofinal}{{04E6}{source}{}{sourcepointer-3}{}}
\newlabel{categories-definition-cofinal@cref}{{[originaldefinition][0][]04E6}{[1][0][]source}{}{}{}}
\makeatother
\crefname{originalsection}{original section}{original sections}
\Crefname{originalsection}{Original section}{Original sections}
\makeatletter
\newlabel{categories-remark-left-localization-morphisms-colimit}{{05Q0}{source}{}{sourcepointer-4}{}}
\newlabel{categories-remark-left-localization-morphisms-colimit@cref}{{[originalsection][0][]05Q0}{[1][0][]source}{}{}{}}
\makeatother
\title{Filtered localization systems of distinguished triangles and cofinality of component functors}
\author{The Stacks Project Authors}
\date{stacks-project:derived.tex:local-snapshot-20261011; GNU FDL 1.2 or later}
\begin{document}
\maketitle
\noindent\textit{Editorial note.} References to standard Categories inputs display their stable Stacks tags. The corresponding permitted inputs are explicitly stated in task metadata.\par\medskip
\noindent\textit{Editorial note.} Triangle compatibility means MS5 (shift stability of S) and MS6 (completion of two S-maps to a third S-map in a morphism of distinguished triangles). These are hypotheses, not unproved solution lemmas.\par\medskip
\noindent\textit{Editorial scope.} One original lemma. Standard calculus of fractions is explicitly allowed; all source-local triangle coherence remains in original proof.\par\medskip
\noindent\textit{Allowed standard prerequisites (editorial).}\begin{itemize}
\item Calculus of fractions for a saturated multiplicative system. Ore completion, common-denominator equalization, filtered categories X/S, and equality in localization after postcomposition by an S-map. Saturation says any map inverted by localization lies in S.. Categories tags05Q0 and05Q9; these are standard calculus-of-fractions inputs, not triangle coherence.
\item Pre-triangulated rotation and compatibility MS5/MS6. S is stable under shift; every commutative square between first two entries of distinguished triangles with both maps in S admits a third S-map completing the triangle morphism.. derived.tex1156-1181
\item Filtered categories and cofinal functors. Filtered means nonempty, common upper targets and equalization of parallel arrows; cofinal means each target receives an arrow and the corresponding comma categories are connected.. Categories tags002V and04E6
\end{itemize}\medskip
\ifdefined\hypertarget\hypertarget{sourcepointer-1}{}\fi
\noindent\textit{Source reference.} Original lemma 05Q9: focused-seeds/ACL/sources/stacks-snapshot-20261011/categories.tex, source lines 4795--4811. Explicit permitted standard input; the target triangle-coherence construction is not supplied by this reference. Display number is stable Stacks tag, not chapter numbering.\par\medskip
\ifdefined\hypertarget\hypertarget{sourcepointer-2}{}\fi
\noindent\textit{Source reference.} Original definition 002V: focused-seeds/ACL/sources/stacks-snapshot-20261011/categories.tex, source lines 2130--2147. Explicit permitted standard input; the target triangle-coherence construction is not supplied by this reference. Display number is stable Stacks tag, not chapter numbering.\par\medskip
\ifdefined\hypertarget\hypertarget{sourcepointer-3}{}\fi
\noindent\textit{Source reference.} Original definition 04E6: focused-seeds/ACL/sources/stacks-snapshot-20261011/categories.tex, source lines 1725--1751. Explicit permitted standard input; the target triangle-coherence construction is not supplied by this reference. Display number is stable Stacks tag, not chapter numbering.\par\medskip
\ifdefined\hypertarget\hypertarget{sourcepointer-4}{}\fi
\noindent\textit{Source reference.} Original section 05Q0: focused-seeds/ACL/sources/stacks-snapshot-20261011/categories.tex, source lines 4362--4399. Explicit permitted standard input; the target triangle-coherence construction is not supplied by this reference. Display number is stable Stacks tag, not chapter numbering.\par\medskip
\begin{lemma}
\label{lemma-limit-triangles}
Let $\mathcal{D}$ be a pre-triangulated category.
Let $S$ be a saturated multiplicative system in $\mathcal{D}$
that is compatible with the triangulated structure.
Let $(X, Y, Z, f, g, h)$ be a distinguished triangle in $\mathcal{D}$.
Consider the category of morphisms of triangles
$$
\mathcal{I} =
\{(s, s', s'') : (X, Y, Z, f, g, h) \to (X', Y', Z', f', g', h')
\mid s, s', s'' \in S\}
$$
Then $\mathcal{I}$ is a filtered category and the functors
$\mathcal{I} \to X/S$, $\mathcal{I} \to Y/S$, and $\mathcal{I} \to Z/S$
are cofinal.
\end{lemma}

\begin{proof}
We strongly suggest the reader skip the proof of this lemma and instead
work it out on a napkin.

\medskip\noindent
The first remark is that using rotation of distinguished triangles (TR2)
gives an equivalence of categories between $\mathcal{I}$ and the
corresponding category for the distinguished triangle
$(Y, Z, X[1], g, h, -f[1])$. Using this we see for example that if
we prove the functor $\mathcal{I} \to X/S$ is cofinal, then
the same thing is true for the functors $\mathcal{I} \to Y/S$ and
$\mathcal{I} \to Z/S$.

\medskip\noindent
Note that if $s : X \to X'$ is a morphism of $S$, then using
MS2 we can find $s' : Y \to Y'$ and $f' : X' \to Y'$ such that
$f' \circ s = s' \circ f$, whereupon we can use MS6 to complete
this into an object of $\mathcal{I}$. Hence the functor
$\mathcal{I} \to X/S$ is surjective on objects. Using rotation as above
this implies the same thing is true for the functors
$\mathcal{I} \to Y/S$ and $\mathcal{I} \to Z/S$.

\medskip\noindent
Suppose given objects $s_1 : X \to X_1 $ and $s_2 : X \to X_2$ in
$X/S$ and a morphism $a : X_1 \to X_2$ in $X/S$. Since $S$ is saturated,
we see that $a \in S$, see
Categories, Lemma \ref{categories-lemma-what-gets-inverted}.
By the argument of the previous paragraph we can complete
$s_1 : X \to X_1$ to an object
$(s_1, s'_1, s''_1) : (X, Y, Z, f, g, h) \to (X_1, Y_1, Z_1, f_1, g_1, h_1)$
in $\mathcal{I}$. Then we can repeat and find
$(a, b, c) : (X_1, Y_1, Z_1, f_1, g_1, h_1) \to (X_2, Y_2, Z_2, f_2, g_2, h_2)$
with $a, b, c \in S$ completing the given $a : X_1 \to X_2$.
But then $(a, b, c)$ is a morphism in $\mathcal{I}$.
In this way we conclude that the functor $\mathcal{I} \to X/S$ is
also surjective on arrows. Using rotation as above,
this implies the same thing is true for the functors
$\mathcal{I} \to Y/S$ and $\mathcal{I} \to Z/S$.

\medskip\noindent
The category $\mathcal{I}$ is nonempty as the identity provides an object.
This proves the condition (1) of the definition of a filtered category, see
Categories, Definition \ref{categories-definition-directed}.

\medskip\noindent
We check condition (2) of
Categories, Definition \ref{categories-definition-directed}
for the category $\mathcal{I}$. Suppose given objects
$(s_1, s'_1, s''_1) : (X, Y, Z, f, g, h) \to (X_1, Y_1, Z_1, f_1, g_1, h_1)$
and
$(s_2, s'_2, s''_2) : (X, Y, Z, f, g, h) \to (X_2, Y_2, Z_2, f_2, g_2, h_2)$
in $\mathcal{I}$. We want to find an object of $\mathcal{I}$
which is the target of an arrow from both
$(X_1, Y_1, Z_1, f_1, g_1, h_1)$ and $(X_2, Y_2, Z_2, f_2, g_2, h_2)$.
By Categories, Remark
\ref{categories-remark-left-localization-morphisms-colimit}
the categories $X/S$, $Y/S$, $Z/S$ are filtered.
Thus we can find $X \to X_3$ in $X/S$ and morphisms
$s : X_2 \to X_3$ and $a : X_1 \to X_3$. By the above we can find a morphism
$(s, s', s'') : (X_2, Y_2, Z_2, f_2, g_2, h_2) \to
(X_3, Y_3, Z_3, f_3, g_3, h_3)$ with $s', s'' \in S$.
After replacing $(X_2, Y_2, Z_2)$ by $(X_3, Y_3, Z_3)$ we may
assume that there exists a morphism $a : X_1 \to X_2$ in $X/S$.
Repeating the argument for $Y$ and $Z$ (by rotating as above)
we may assume there is a morphism
$a : X_1 \to X_2$ in $X/S$,
$b : Y_1 \to Y_2$ in $Y/S$, and
$c : Z_1 \to Z_2$ in $Z/S$.
However, these morphisms do not necessarily give rise to a morphism of
distinguished triangles. On the other hand, the necessary diagrams
do commute in $S^{-1}\mathcal{D}$. Hence we see (for example) that
there exists a morphism $s'_2 : Y_2 \to Y_3$ in $S$ such that
$s'_2 \circ f_2 \circ a = s'_2 \circ b \circ f_1$. Another replacement
of $(X_2, Y_2, Z_2)$ as above then gets us to the situation where
$f_2 \circ a = b \circ f_1$. Rotating and applying the same argument
two more times we see that we may assume $(a, b, c)$ is a morphism
of triangles. This proves condition (2).

\medskip\noindent
Next we check condition (3) of
Categories, Definition \ref{categories-definition-directed}.
Suppose $(s_1, s_1', s_1'') : (X, Y, Z) \to (X_1, Y_1, Z_1)$ and
$(s_2, s_2', s_2'') : (X, Y, Z) \to (X_2, Y_2, Z_2)$
are objects of $\mathcal{I}$, and suppose $(a, b, c), (a', b', c')$
are two morphisms between them. Since $a \circ s_1 = a' \circ s_1$
there exists a morphism $s_3 : X_2 \to X_3$ such that
$s_3 \circ a = s_3 \circ a'$. Using the surjectivity statement
we can complete this to a morphism of triangles
$(s_3, s_3', s_3'') : (X_2, Y_2, Z_2) \to (X_3, Y_3, Z_3)$
with $s_3, s_3', s_3'' \in S$. Thus
$(s_3 \circ s_2, s_3' \circ s_2', s_3'' \circ s_2'') :
(X, Y, Z) \to (X_3, Y_3, Z_3)$ is also an object of $\mathcal{I}$
and after composing the maps $(a, b, c), (a', b', c')$ with
$(s_3, s_3', s_3'')$ we obtain $a = a'$. By rotating we may do the
same to get $b = b'$ and $c = c'$.

\medskip\noindent
Finally, we check that $\mathcal{I} \to X/S$ is cofinal, see
Categories, Definition \ref{categories-definition-cofinal}.
The first condition is true as the functor is surjective.
Suppose that we have an object $s : X \to X'$ in $X/S$ and
two objects
$(s_1, s'_1, s''_1) : (X, Y, Z, f, g, h) \to (X_1, Y_1, Z_1, f_1, g_1, h_1)$
and
$(s_2, s'_2, s''_2) : (X, Y, Z, f, g, h) \to (X_2, Y_2, Z_2, f_2, g_2, h_2)$
in $\mathcal{I}$ as well as morphisms $t_1 : X' \to X_1$ and
$t_2 : X' \to X_2$ in $X/S$. By property (2) of $\mathcal{I}$
proved above we can find morphisms
$(s_3, s'_3, s''_3) : (X_1, Y_1, Z_1, f_1, g_1, h_1) \to
(X_3, Y_3, Z_3, f_3, g_3, h_3)$
and
$(s_4, s'_4, s''_4) : (X_2, Y_2, Z_2, f_2, g_2, h_2) \to
(X_3, Y_3, Z_3, f_3, g_3, h_3)$ in $\mathcal{I}$.
We would be done if the compositions
$X' \to X_1 \to X_3$ and $X' \to X_2 \to X_3$ were equal
(see displayed equation in 
Categories, Definition \ref{categories-definition-cofinal}).
If not, then, because $X/S$ is filtered, we can choose
a morphism $X_3 \to X_4$ in $X/S$ such that the compositions
$X' \to X_1 \to X_3 \to X_4$ and $X' \to X_2 \to X_3 \to X_4$ are equal.
Then we finally complete $X_3 \to X_4$ to a morphism
$(X_3, Y_3, Z_3) \to (X_4, Y_4, Z_4)$ in $\mathcal{I}$
and compose with that morphism to see that the result is true.
\end{proof}

\end{document}
